\subsection{Macaulay 模的支撑集}

命题 2. — 设 A 为 Noether 环，M 为有限生成 Macaulay A-模。

\begin{enumerate}
\def\labelenumi{\alph{enumi})}
\item
  A-模 M 没有嵌入相伴素理想。 \(^{1}\)
\item
  设 X 为 Supp(M) 的闭不可约子集，Y 为 X 的闭子集。有
\end{enumerate}

\[
\operatorname{codim} (Y, X) + \operatorname{codim} (X, \operatorname{Supp} (M)) = \operatorname{codim} (Y, \operatorname{Supp} (M)).
\]

\begin{enumerate}
\def\labelenumi{\alph{enumi})}
\setcounter{enumi}{2}
\item
  拓扑空间 Supp(M) 是链式的 (VIII, § 1, 第 2 节, 定义 4)。
\item
  设 \(X_{1}\) 与 \(X_{2}\) 为 Supp(M) 的不可约分支，Y 为 \(X_{1} \cap X_{2}\) 的闭子集。有 codim(Y, \(X_{1}\) ) = codim(Y, \(X_{2}\) )。
\item
  设 \(\mathfrak{p} \in \operatorname{Ass}(M)\) 。有 \(\operatorname{prof}_{\mathrm{A}}(\mathfrak{p}; M) = 0\) （ \(\S\) 1, \(n^{\circ}\) 1, 注 2），从而 \(\dim(M_{\mathfrak{p}}) = 0\) （ \(n^{\circ}\) 1, 命题 1 的推论），这蕴含 \(\mathfrak{p}\) 是 \(\operatorname{Supp}(M)\) 的极小元。
\item
  先设 Y 不可约。设 p 与 q 为 Supp(M) 的素理想，使得 \(Y = V(\mathfrak{q})\) 且 \(X = V(\mathfrak{p})\) 。由命题 1 我们有
\end{enumerate}

\[
\begin{array}{r l} \operatorname{codim} (Y, X) & = \dim (A _ {q} / p A _ {q}) = \dim (M _ {q}) - \dim (M _ {p}) \\ & = \operatorname{codim} (Y, \operatorname{Supp} (M)) - \operatorname{codim} (X, \operatorname{Supp} (M)). \end{array}
\]

一般情形由 VIII, § 1, 第 2 节, 注 3 得出。

\begin{enumerate}
\def\labelenumi{\alph{enumi})}
\setcounter{enumi}{2}
\tightlist
\item
  设 X、Y、Z 为 Supp(M) 的闭不可约子集且 \(Z \subset Y \subset X\) 。这些子集在 Supp(M) 中的余维数均有限 (VIII, § 1, 第 4 节, 命题 9 及 § 3, 第 1 节, 命题 2 的推论 1)。于是由 b) 推出等式
\end{enumerate}

\[
\operatorname{codim} (Z, Y) + \operatorname{codim} (Y, X) = \operatorname{codim} (Z, X)
\]

它蕴含 c) (VIII, § 1, 第 2 节, 命题 4)。

\begin{enumerate}
\def\labelenumi{\alph{enumi})}
\setcounter{enumi}{3}
\tightlist
\item
  由 b)，有 \(\text{codim}(Y, X_1) = \text{codim}(Y, \text{Supp}(M)) = \text{codim}(Y, X_2)\) 。
\end{enumerate}

特别地，若存在支撑集等于 Spec(A) 的有限生成 Macaulay A-模 M，则环 A 是链式的，从而 A 的每个分式环或每个商环都是链式的 (VIII, § 1, 第 3 节, 注 2)。

注 1. — 在命题 2 的假设下，可能发生这样的情形：Supp(M) 的两个不可约分支 \(X_{1}\) 与 \(X_{2}\) 的交 Y 退化为一点，且有 \(\dim X_{1} \neq \dim X_{2}\) 及 \(\dim X_{2} \neq \text{codim}(Y, X_{2})\) （参见习题 4）。然而，当环 A 是局部环时这不可能发生，如下面的推论所示。

推论. — 设 A 为 Noether 局部环，M 为非零的有限生成 Macaulay A-模。

\begin{enumerate}
\def\labelenumi{\alph{enumi})}
\item
  Supp(M) 的闭不可约子集的所有极大链长度都等于 dim(M)。
\item
  对 Supp(M) 的每个闭子集 X，有
\end{enumerate}

\[
\operatorname{codim} (\mathrm{X}, \operatorname{Supp} (\mathrm{M})) = \dim (\operatorname{Supp} (\mathrm{M})) - \dim (\mathrm{X}).
\]

\begin{enumerate}
\def\labelenumi{\alph{enumi})}
\setcounter{enumi}{2}
\item
  Supp(M) 的所有不可约分支维数相同。
\item
  对 A 的每个理想 J，有
\end{enumerate}

\[
\operatorname{prof} _ {\mathrm{A}} (\mathrm{J}; \mathrm{M}) = \dim (\mathrm{M}) - \dim (\mathrm{M} / \mathrm{JM}).
\]

\begin{enumerate}
\def\labelenumi{\alph{enumi})}
\tightlist
\item
  Supp(M) 的闭不可约子集的一个极大链，其最小元为 \(\{\mathfrak{m}_{\mathrm{A}}\}\) ，其最大元为不可约分支 X，X 属于
\end{enumerate}

Supp(M)。它的长度等于 \(\{m_{A}\}\) 在 X 中的余维数（命题 2, c)）；由命题 2, b) 应用于闭子集 \(\{m_{A}\}\subset X\) ，它等于 \(\text{codim}(\{\mathfrak{m}_{\Lambda}\},\text{Supp}(M))\) ，即 \(\dim(M)\) 。

\begin{enumerate}
\def\labelenumi{\alph{enumi})}
\setcounter{enumi}{1}
\item
  当子集 X 不可约时这是 a) 的推论 (VIII, § 1, 第 2 节, 命题 5)；一般情形由 VIII, § 1, 第 1 节, 命题 1 及 § 1, 第 2 节, 注 4 得出。
\item
  这是 b) 的推论。
\item
  有 \(\operatorname{prof}_{\mathrm{A}}(\mathrm{J};\mathrm{M}) = \operatorname{codim}(\operatorname{Supp}(\mathrm{M})\cap \mathrm{V}(\mathrm{J}),\operatorname{Supp}(\mathrm{M}))\) ，由第 1 节的命题 1 的推论。于是只需取 \(\mathrm{X} = \operatorname{Supp}(\mathrm{M})\cap \mathrm{V}(\mathrm{J}) = \operatorname{Supp}(\mathrm{M} / \mathrm{JM})\) 应用 b) (II, § 4, 第 4 节, 命题 18 的推论)。
\end{enumerate}

注 2.--- 设 M 为有限生成 Macaulay A-模，p 为 Supp(M) 的元素。由 § 1, 第 4 节的定理 2，有 prof\_A(p; M) \textless{} +∞，且存在长度为 prof\_A(p; M) 的、由 p 中元素组成的 M-正则序列。用 J 表示由这样一个序列生成的 A 的理想；则 A-模 M/JM 是 Macaulay 的（第 1 节, 例 3），且 p 是其支撑集的极小元。事实上，p 包含 J，从而属于 M/JM 的支撑集 (II, § 4, 第 4 节, 命题 18 的推论)；由 § 1, 第 4 节的定理 2 的推论 1，理想 p 包含于 Ass(M/JM) 的某个元素，但与有限生成 Macaulay 模相伴的每个素理想都是其支撑集的极小元（命题 2）。
% semantic-fragments: legacy-input-seam
\relax{}
